% Options for packages loaded elsewhere
\PassOptionsToPackage{unicode}{hyperref}
\PassOptionsToPackage{hyphens}{url}
\PassOptionsToPackage{dvipsnames,svgnames,x11names}{xcolor}
\documentclass[
  10pt,
  fontset=fandol]{ctexart}
\usepackage{xcolor}
\usepackage[margin=16mm]{geometry}
\usepackage{amsmath,amssymb}
\setcounter{secnumdepth}{-\maxdimen} % remove section numbering
\usepackage{iftex}
\ifPDFTeX
  \usepackage[T1]{fontenc}
  \usepackage[utf8]{inputenc}
  \usepackage{textcomp} % provide euro and other symbols
\else % if luatex or xetex
  \usepackage{unicode-math} % this also loads fontspec
  \defaultfontfeatures{Scale=MatchLowercase}
  \defaultfontfeatures[\rmfamily]{Ligatures=TeX,Scale=1}
\fi
\usepackage{lmodern}
\ifPDFTeX\else
  % xetex/luatex font selection
\fi
% Use upquote if available, for straight quotes in verbatim environments
\IfFileExists{upquote.sty}{\usepackage{upquote}}{}
\IfFileExists{microtype.sty}{% use microtype if available
  \usepackage[]{microtype}
  \UseMicrotypeSet[protrusion]{basicmath} % disable protrusion for tt fonts
}{}
\makeatletter
\@ifundefined{KOMAClassName}{% if non-KOMA class
  \IfFileExists{parskip.sty}{%
    \usepackage{parskip}
  }{% else
    \setlength{\parindent}{0pt}
    \setlength{\parskip}{6pt plus 2pt minus 1pt}}
}{% if KOMA class
  \KOMAoptions{parskip=half}}
\makeatother
\usepackage{longtable,booktabs,array}
\newcounter{none} % for unnumbered tables
\usepackage{calc} % for calculating minipage widths
% Correct order of tables after \paragraph or \subparagraph
\usepackage{etoolbox}
\makeatletter
\patchcmd\longtable{\par}{\if@noskipsec\mbox{}\fi\par}{}{}
\makeatother
% Allow footnotes in longtable head/foot
\IfFileExists{footnotehyper.sty}{\usepackage{footnotehyper}}{\usepackage{footnote}}
\makesavenoteenv{longtable}
\setlength{\emergencystretch}{3em} % prevent overfull lines
\providecommand{\tightlist}{%
  \setlength{\itemsep}{0pt}\setlength{\parskip}{0pt}}
\usepackage{amsmath,amssymb,mathtools,needspace}
\xeCJKDeclareCharClass{CJK}{"2032,"0400->"04FF,"2160->"216B,"2460->"2473,"25A0,"25C7,"25CB,"25CF}
\IfFontExistsTF{Arial Unicode MS}{\xeCJKsetup{AutoFallBack=true}\setCJKfallbackfamilyfont{\CJKrmdefault}{Arial Unicode MS}}{}
\setcounter{secnumdepth}{0}
\setlength{\emergencystretch}{2em}
\allowdisplaybreaks
\usepackage{bookmark}
\IfFileExists{xurl.sty}{\usepackage{xurl}}{} % add URL line breaks if available
\urlstyle{same}
\hypersetup{
  pdftitle={改进的 Euler 格式},
  colorlinks=true,
  linkcolor={Maroon},
  filecolor={Maroon},
  citecolor={Blue},
  urlcolor={Blue},
  pdfcreator={LaTeX via pandoc}}

\title{改进的 Euler 格式}
\author{}
\date{}

\begin{document}
\maketitle

\subsubsection{5.2.4 改进的 Euler 格式}\label{ux6539ux8fdbux7684-euler-ux683cux5f0f}

我们看到，梯形方法虽然提高了精度，但其算法复杂，在应用迭代公式 (5.2.9) 进行实际计算时，每迭代一次，都要重新计算函数 \(f\) 的值，而迭代又要反复进行若干次，计算量很大，而且往往难以预测。为了控制计算量，通常希望只迭代一两次就转入下一步计算，从而简化算法。

具体地说，先用 Euler 格式求得一个初步的近似值 \(\bar y_{n+1}\)，称之为\textbf{预测值}。预测值 \(\bar y_{n+1}\) 的精度可能很差，再用梯形公式 (5.2.8) 将它校正一次，即按式 (5.2.9) 迭代一次，得 \(y_{n+1}\)，这个结果称为\textbf{校正值}。而这样建立的预测-校正系统通常称为\textbf{改进的 Euler 格式}：

预测

\[
\bar y_{n+1}=y_n+hf(x_n,y_n),\tag{5.2.10}
\]

校正

\[
y_{n+1}=y_n+\frac{h}{2}[f(x_n,y_n)+f(x_{n+1},\bar y_{n+1})].\tag{5.2.11}
\]

这一计算格式亦可表示为

\[
y_{n+1}=y_n+\frac{h}{2}[f(x_n,y_n)+f(x_n+h,y_n+hf(x_n,y_n))],\tag{5.2.12}
\]

或表示为下列平均化形式：

\[
\begin{cases}
y_p=y_n+hf(x_n,y_n),\\
y_c=y_n+hf(x_{n+1},y_p),\\
y_{n+1}=\dfrac{1}{2}(y_p+y_c).
\end{cases}
\]

\href{../../source-images/pdf-124.jpeg}{原页}

\textbf{例 5.2} 用改进的 Euler 方法求解初值问题 (5.2.2)。

\textbf{解} 改进的 Euler 格式为

\[
\begin{cases}
y_p=y_n+h\left(y_n-\dfrac{2x_n}{y_n}\right),\\
y_c=y_n+h\left(y_p-\dfrac{2x_{n+1}}{y_p}\right),\\
y_{n+1}=\dfrac{1}{2}(y_p+y_c).
\end{cases}
\]

仍取 \(h=0.1\)，计算结果如表 5.2 所示。与例 5.1 Euler 方法的计算结果比较，改进的 Euler 方法明显地改善了精度。

\textbf{表 5.2}

{\def\LTcaptype{none} % do not increment counter
\begin{longtable}[]{@{}llllll@{}}
\toprule\noalign{}
\(x_n\) & \(y_n\) & \(y(x_n)\) & \(x_n\) & \(y_n\) & \(y(x_n)\) \\
\midrule\noalign{}
\endhead
\bottomrule\noalign{}
\endlastfoot
0.1 & 1.0959 & 1.0954 & 0.6 & 1.4860 & 1.4832 \\
0.2 & 1.1841 & 1.1832 & 0.7 & 1.5525 & 1.5492 \\
0.3 & 1.2662 & 1.2649 & 0.8 & 1.6153 & 1.6125 \\
0.4 & 1.3434 & 1.3416 & 0.9 & 1.6782 & 1.6733 \\
0.5 & 1.4164 & 1.4142 & 1.0 & 1.7379 & 1.7321 \\
\end{longtable}
}

\begin{center}\rule{0.5\linewidth}{0.5pt}\end{center}

\subsection{相关原页校注（agent 补充）}\label{ux76f8ux5173ux539fux9875ux6821ux6ce8agent-ux8865ux5145}

以下校注来自本节覆盖的原页，可能也涉及同页相邻小节；原文照录与校正说明分开保留。

\subsubsection{原书 PDF 页 124 的校注}\label{ux539fux4e66-pdf-ux9875-124-ux7684ux6821ux6ce8}

\href{../pages/pdf-124.md}{完整原页转录} · \href{../../source-images/pdf-124.jpeg}{原页图像}

校注记录：ch05a-E001, ch05a-E002。

\begin{quote}
\textbf{校注（agent 补充，原书疑误）}：表 5.2 在 \(x_n=0.8\) 的 \(y_n\) 原印为 \(1.6153\)，已按原图保留。按例 5.2 所列递推式从 \(y_0=1\)、\(h=0.1\) 计算，得 \(y_7=1.5525140913\cdots\)、\(y_8=1.6164747827\cdots\)，保留四位小数应为 \(1.6165\)。这是原书数值疑误，非 OCR 修正；其余各行按同一计算与表中四位小数相符。

\textbf{校注（agent 补充，原书下标错误）}：本页``两步格式''说明中的``更前面一步的信息 \(y_{n+1}\)''原图确实印为 \(n+1\)；按紧邻的式 (5.2.13)，此处应指 \(y_{n-1}\)。正文保留原印下标。
\end{quote}

\subsection{本节覆盖原页的校注记录}\label{ux672cux8282ux8986ux76d6ux539fux9875ux7684ux6821ux6ce8ux8bb0ux5f55}

下列记录按原页关联，含同页相邻小节的校注；计算时同时读取上文校注和完整原页。

\begin{itemize}
\tightlist
\item
  \texttt{ch05a-E001}：原书 PDF 页 124
\item
  \texttt{ch05a-E002}：原书 PDF 页 124
\end{itemize}

\end{document}
